\documentclass[11pt]{article}
\usepackage{fontspec}
\usepackage[french]{babel}
\usepackage{amsmath}
\usepackage{amssymb}
\usepackage{amsthm}
\usepackage{mathrsfs}
\usepackage{tikz-cd}
\usepackage{enumitem}

\newcommand{\printedpage}[1]{\marginpar{\footnotesize p.~#1}}
\newcommand{\msp}[1]{\marginpar{\footnotesize [#1]}}
\newcommand{\defeq}{\overset{\text{déf}}{=}}
\newcommand{\fnm}[1]{\textsuperscript{(}\footnotemark[#1]\textsuperscript{)}}
\newcommand{\cl}{^{\tilde{}}}% tilde printed as exponent after a symbol or a parenthesis

\begin{document}

\printedpage{39}
% head: A. Grothendieck : Dérivateurs, Ch. XVII, Édition des Universités de Montpellier II et Paris VII
qui se décompose en deux
\[
\left\{
\begin{array}{rcl}
v_{!}(TC) & \subset & TC' \\
T\tilde{C}' & \subset & (v_{!}(TC))\cl.
\end{array}
\right.
\tag*{(4.7.3)}
\]
La moins délicate des deux est la \msp{page 80} première, $v_{!}(TC) \subset TC'$, qui revient, en vertu des définitions, impliquant
\[
v_{!}(TC) \subset v_{!}(C) \subset C' \, ,
\]
à
\[
v_{!}(TC) \;\subset\; W' \; ,
\]
i.e. par définition de $W'$, à
\[
v^{*}v_{!}(TC) \;\subset\; TC \; ,
\tag*{(4.7.4)}
\]
ce qui est une hypothèse de type (4.6.1) $3^{\circ}$), à joindre aux hypothèses (4.6.1).

La relation hypothétique
\[
\underbrace{T\tilde{C}'}_{=\,\tilde{C}'\cap W'} \;\subset\; (v_{!}(TC))\cl, \quad \text{ou seulement} \quad TC' \subset (v_{!}(TC))\cl,
\tag*{$(*)$}
\]
m'apparaît de nature nettement plus délicate. Comme $T\tilde{C}' \supset TC' \supset v_{!}(TC)$, elle équivaut à l'\textit{égalité} de $(T\tilde{C}')_{*}$ resp. $(TC')_{*}$ et $(v_{!}(TC))_{*} = (v^{*})^{-1}(TC)_{*}$. Posons
\[
\left\{
\begin{array}{lcl}
(TC)_{*} & \defeq & F \\
(TC')_{*} & \defeq & F' \;\supset\; (T\tilde{C}')_{*} \;=\; F'_{\sim} \; .
\end{array}
\right.
\tag*{(4.7.5)}
\]
On a donc, en vertu de l'inclusion
\[
v_{!}(TC) \;\subset\; TC' \; ( \;\subset\; T\tilde{C}' \; ) \; ,
\]
la relation
\[
F'_{\sim} \;\subset\; F' \;\subset\; (v^{*})^{-1}(F) \; ,
\tag*{(4.7.6)}
\]
et la question est donc si on a égalité $F' = (v^{*})^{-1}(F)$, voire $F'_{\sim} = (v^{*})^{-1}(F)$ (d'où $F'_{\sim} = F'$).

\msp{page 81}

En fait, on montre assez facilement qu'on a $F'_{\sim} = F'$, i.e. $T\tilde{C}' \defeq \tilde{C}' \cap W' \subset (C' \cap W')\cl$ -- c'est une histoire générale de catégories à cofibrations spéciales \textsuperscript{(}\footnote[21]{cf. prop. 5.1 ci-dessous.}\textsuperscript{)}. Mais je ne vois guère comment prouver (sous des hypothèses convenables) que $F' = (v^{*})^{-1}(F)$. Il me faudra revenir là-dessus au besoin \textsuperscript{(}\footnote[22]{voir cependant prop. 5.8 ci-dessous.}\textsuperscript{)}. Pour l'instant, je me bornerai à récolter le fruit de 4.6 dans un cas particulier :

\textbf{Proposition 4.8.} \textit{Soit} $(\mathcal{M}, W, C)$ \textit{une catégorie à cofibrations spéciale} (déf. 4.4, page 94)\textit{. Soit $I$ objet de} Cat\textit{, $I_0$ la catégorie discrète associée, considérons $v : I_0 \longrightarrow I$ et} % [sic] page 94 -> page 74, apparently: the same déf. 4.4 is cited as "page 74" on printed p. 42
% foot: Version 9 février 2019
\printedpage{40}
% head: A. Grothendieck : Dérivateurs, Ch. XVII, Édition des Universités de Montpellier II et Paris VII
\[
\begin{tikzcd}
\mathcal{M}(I_0) \arrow[r, shift left, "v_{!}"] & \mathcal{M}(I) \arrow[l, shift left, "v^{*}"]
\end{tikzcd} \; ,
\]
\textit{enfin munissons $\mathcal{M}(I_0) = \mathcal{M}^{I_0}$ de la structure de catégorie à cofibrations,} produit \textit{de celles des facteurs, i.e. définissons $W(I_0)$, $C(I_0)$ “fibre à fibre”. (Cela coïncide avec la définition générale donnée au paragraphe} 2\textit{, pour $I_0$ muni de la structure d'ordre discret, et fait de $\mathcal{M}(I_0)$ une catégorie à cofibrations spéciale.) Alors le couple $(v_{!}, v^{*})$ satisfait} \msp{page 82} \textit{aux conditions de la proposition} 4.6\textit{, ainsi qu'à la condition supplémentaire à} (4.6.1)\textit{,}
\[
4^{\circ}) \qquad v^{*}(v_{!}(TC(I))) \;\subset\; TC(I) \; .
\tag*{(4.8.1)}
\]
\textit{Donc on trouve une structure de catégorie à cofibrations spéciale sur $\mathcal{M}(I)$, en posant}
\[
W(I) \;=\; (v^{*})^{-1}(W(I_0))
\tag*{(4.8.2)}
\]
\textit{(c'est en fait équivalent à la définition habituelle de $W(I)$) et}
\[
C(I) \quad = \quad
\begin{array}[t]{l}
\textit{clôture de } v_{!}(C(I_0)) \textit{ pour les opérations de cochangement de}\\
\textit{base et composition transfinie.}
\end{array}
\tag*{(4.8.3)}
\]
\textit{On a de plus sous ces conditions :}
\[
\underbrace{TF(I)}_{\defeq\,(C(I))_{*}} \;=\; (v^{*})^{-1}( \; \underbrace{TF(I_0)}_{\substack{\defeq\,(C_{I_0})_{*},\\ \textit{se définit aussi}\\ \textit{fibre par fibre}}} \; ) \; ,
\tag*{(4.8.8)}
\]
\[
\left\{
\begin{array}{rcl}
v_{!}( \, \overbrace{(TC)(I_0)}^{C(I_0)\cap W(I_0)} \, ) & \subset & TC(I) \qquad\qquad\qquad \text{\fnm{23}} \\
F(I) \defeq (TC(I))_{*} & \subset & (v^{*})^{-1}(F(I_0)).
\end{array}
\right.
\tag*{(4.8.9)}
\]
\footnotetext[23]{Ces inclusions devraient être des égalités, mais je ne sais pas le prouver sans faire des hypothèses plus fortes, du type \textit{catégorie de Quillen spéciale.}}

\textsc{Démonstration}. On sait déjà que $\mathcal{M}(I)$ est stable par petites $\varinjlim$ , et que tout élément est accessible, comme il résulte aussitôt des propriétés idoines de $\mathcal{M}$. Il est clair aussi que $v^{*}$ commute aux petites $\varinjlim$ . Il reste à prouver (4.6.1. $3^{\circ}$), et $4^{\circ}$ dans (4.8.1), i.e. \msp{page 83} que $v^{*}v_{!} : \mathcal{M}(I_0) \longrightarrow \mathcal{M}(I_0)$ applique $C$ dans $C$, $TC$ dans $TC$. Or on a
\[
v_{!}(F_0)(i) \;=\; \varinjlim_{(\alpha,j)\in i\backslash I_0} F_0(j) \qquad\quad \text{\fnm{24}}
\]
\footnotetext[24]{où $j \in \mathrm{Ob}\, I_0 = \mathrm{Ob}\, I$, et $\alpha : i \longrightarrow v(j) = j$.}
et $i\backslash I_0$ est la catégorie discrète au dessus de $I_0$, dont la fibre en $j \in \mathrm{Ob}\, I_0 \;=\; \mathrm{Ob}\, I$ est $\mathrm{Hom}_I(i,j)$ :
\[
i\backslash I_0 \;\;=\;\; \coprod_{j\in \mathrm{Ob}\, I} \text{(catégorie discrète associée à $\mathrm{Hom}_I(i,j)$)} \; ,
\]
% foot: Version 9 février 2019
\printedpage{41}
% head: A. Grothendieck : Dérivateurs, Ch. XVII, Édition des Universités de Montpellier II et Paris VII
donc
\[
v_{!}(F_0)(i) \;=\; \coprod_{j\in \mathrm{Ob}\, I = \mathrm{Ob}\, I_0} \underbrace{F_0(j)^{\mathrm{Hom}_I(i,j)}}_{\substack{\text{somme directe}\\ \text{de } \mathrm{Hom}_I(i,j)\\ \text{copies de } F_0(j)}} \, ,
\]
donc on trouve
\[
v^{*}v_{!}F_0( \, \underbrace{\; i \;}_{\in \mathrm{Ob}\, I_0 \,=\, \mathrm{Ob}\, I} \, ) \;=\; v_{!}(F_0)(i) \;=\; \text{\tikz[baseline=0pt]\draw[->] (0,0) .. controls (0.4,0) .. (0.5,0.45);} \quad .
\]
% the curved arrow points up to the underbraced expression F_0(j)^{Hom_I(i,j)} of the preceding display
Comme $C$ et $TC$ sont stables par sommes directes quelconques (l'étant par sommes directes finies, à cause des conditions de stabilité finies sur $C$ et $TC$, et aussi par somme infinies à cause de l'axiome de stabilité transfinie \dots\textsuperscript{(}\footnote[25]{mérite un lemme général \dots}\textsuperscript{)}), il s'ensuit bien que $v^{*}v_{!}$ applique $C(I_0)$ et $TC(I_0)$ en eux-mêmes. % [sic] somme infinies -> sommes infinies, agreement with the plural adjective

Ainsi, la première assertion de 4.8 résulte de 4.6. De plus, (4.8.8) et (4.8.9) ont été vus dans 4.7, dans les conditions générales de 4.6 et en supposant de plus le $4^{\circ}$ de (4.7.4).

\msp{page 84}

\section*{5\quad Relation des catégories à cofibrations spéciales avec les catégories de Quillen spéciales}

Soit $(\mathcal{M}, W, C)$ une catégorie à cofibrations spéciale. J'introduis
\[
TF \;=\; C_{*}, \qquad\quad TC \;=\; C \cap W, \qquad\quad F \;=\; (TC)_{*} \; ,
\tag*{(5.1)}
\]
donc on a
\[
\begin{tikzcd}[row sep=large, column sep=large]
TC \arrow[r, hook] \arrow[d, leftrightarrow] & C \arrow[d, leftrightarrow, "\shortmid"{marking, sloped}] \\
F & TF \arrow[l, hook']
\end{tikzcd} \; ,
\tag*{(5.2)}
\]
$(C, TF)$ étant un précouple de Quillen clos à droite. On sait alors (XV 4.2 cor. 1, page 31) que $\tilde{C}$ est formé des facteurs directs d'éléments de $C$. Posons
\[
T\tilde{C} \;=\; \tilde{C} \cap W \; ,
\tag*{(5.3)}
\]
d'où
\[
TC \;\subset\; T\tilde{C} \; .
\tag*{(5.4)}
\]

\textbf{Proposition 5.1} \textit{Soit $(\mathcal{M}, W, C)$ une catégorie à cofibrations spéciale. On a $T\tilde{C} \subset (TC)\cl$ (et même $T\tilde{C} \subset (TC)^{\natural}$ \textsuperscript{(}\footnote[26]{\textbf{N.B.} Pour une partie $A$ de $\mathrm{Fl}(\mathcal{M})$, $A^{\natural}$ désigne l'ensemble des facteurs directs d'éléments de $A$.}\textsuperscript{)}), ou ce qui revient au même en vertu de (5.4), $TC$ et $T\tilde{C}$ ont même $Q$-clôture à gauche :}
\[
T\tilde{C} = (T\tilde{C})\cl = (TC)\cl \qquad\quad \textit{(et même } T\tilde{C} \subset (TC)^{\natural}\textit{),}
\tag*{(5.5)}
\]
% foot: Version 9 février 2019
\printedpage{42}
% head: A. Grothendieck : Dérivateurs, Ch. XVII, Édition des Universités de Montpellier II et Paris VII
\textit{d'où}
\[
(\tilde{C})_{*} = C_{*} \, (= TF), \qquad\quad (T\tilde{C})_{*} = (TC)_{*} \, (= F) \qquad\quad \text{\fnm{27}}.
\tag*{(5.6)}
\]
\footnotetext[27]{Ces relations ne dépendent pas de l'hypothèse $W$ stable par facteurs directs.}
\textit{De plus, si $W$ stable par facteurs directs (ou si seulement tout facteur direct d'une flèche de $TC$ est dans $W$), et si de plus $W$ est stable par composition transfinie, alors $(\mathcal{M}, W, \tilde{C})$ est une catégorie à cofibrations spéciale. De plus, $T\tilde{C}$ est $Q$-saturé à gauche, donc} (par 5.5)
\[
T\tilde{C} \;=\; (TC)^{\natural} \; .
\tag*{(5.7)}
\]

\msp{page 85}

\textsc{Démonstration}. Prouvons
\[
T\tilde{C} \;\subset\; (TC)\cl \, ,
\]
et plus précisément, que tout élément de $T\tilde{C}$ est facteur direct d'un élément de $TC$ :
\[
T\tilde{C} = (TC)^{\natural} \; .
\]
Soit donc $f \in T\tilde{C} \defeq \tilde{C}\cap W$, considérons leur factorisation de $f$ en $pi$, avec $i \in C$, $p \in TF$. Comme $TF \subset W$, on a $p \in W$, donc $i \in W$, donc $i \in C \cap W \defeq TC$. Comme $f \in C$, $p \in TF$, il existe $q : Y \longrightarrow X'$ tel que $pq = \mathrm{id}_Y$, $qf = i$,
\[
\begin{tikzcd}[column sep=small, row sep=small]
X \arrow[rr, "f\in C\cap W"] \arrow[dr, "i\in C"'] & & Y \arrow[dl, dotted, shift right, "q"'] \\
& X' \arrow[ur, shift right, "p\in TF"'] &
\end{tikzcd}
\]
donc $f$ apparaît comme facteur direct de $i \in TC$, q.e.d.

Prouvons que $(\mathcal{M}, W, \tilde{C})$ est une catégorie à cofibrations spéciale (déf. 4.4, page 74). Condition a) sur $\mathcal{M}$ acquise, condition b), la partie sur $W$, aussi. La stabilité de $\tilde{C}$ par cochangement de base et par composition OK, aussi (pour b$'$) la stabilité par composition transfinie. Stabilité de $T\tilde{C} = \tilde{C} \cap W$ par cochangement de base : il reste à prouver que si on a un carré cocartésien
\[
\begin{tikzcd}[row sep=large, column sep=large]
X \arrow[r, "i"] \arrow[d] \arrow[dr, phantom, "\text{cocart.}"] & Y \arrow[d] \\
X' \arrow[r, "i'"] & Y'
\end{tikzcd} \; ,
\]
si $i \in \tilde{C}\cap W$, alors $i' \in W$ (car on sait déjà que $i' \in \tilde{C}$). Mais la démonstration précédente montre que c'est facteur direct \textit{dans} $X\backslash\mathcal{M}$ d'une flèche dans $TC$, donc $i'$ est facteur direct d'une flèche de $TC \subset W$. Donc si $W$ est stable par facteurs directs, $i' \in W$.
% foot: Version 9 février 2019

\printedpage{43}
% head: A. Grothendieck : Dérivateurs, Ch. XVII, Édition des Universités de Montpellier II et Paris VII
\msp{page 86}

Cela achève de prouver b). Pour b$'$), il reste à voir que $T\tilde{C}$ est stable par composition transfinie. Je ne vois pas d'autre façon de le voir, qu'en supposant que $W$ lui-même est stable par composition transfinie (ce qui paraît à peine moins exigeant que la stabilité de $W$ par $\varinjlim$ \texttt{[filtrantes]}). Cela achève de prouver b, b$'$). La condition c) sur $C$ implique la même condition sur $\tilde{C}$. Enfin, comme $\tilde{C}_{*} = C_{*}$, la condition d) est la même pour $C$ et $\tilde{C}$. Cela achève de prouver que $(\mathcal{M}, W, \tilde{C})$ est une catégorie à cofibrations spéciale. (\textbf{N.B.} on avait besoin de deux hypothèses supplémentaires sur $W$, stabilité par facteurs directs et par composition transfinie, pour établir que $T\tilde{C}$ est stable par cochangement de base et par composition transfinie). Le dernière assertion de 5.1 et la relation (5.7) qui en résulte, revient à dire que $T\tilde{C} = \tilde{C} \cap W$ est stable par facteurs directs, ce qui est évident par la même propriété de $\tilde{C}$, et l'hypothèse faite sur $W$ de stabilité par facteurs directs (ou seulement que les facteurs directs d'éléments de $T\tilde{C}$ sont dans $W$, donc dans $T\tilde{C}$). % [sic] Le dernière -> La dernière, gender agreement with "assertion"

\msp{page 87}

\textbf{Proposition 5.2.} \textit{Soit $(\mathcal{M}, W, C)$ une catégorie à cofibrations spéciale. Alors avec les notations de} (5.1)\textit{, on a}
\[
W \;=\; TF \circ TC,
\tag*{(5.8)}
\]
\[
TF \;=\; F \cap W.
\tag*{(5.9)}
\]

\textsc{Démonstration}. Soit $f \in W$, écrivons $f = pi$, avec $i \in C$, $p \in TF$, donc $p \in W$, donc $i \in W$, donc $i \in C \cap W = TC$, d'où (5.8).

Prouvons (5.9). On a $TF \subset F \cap W$, prouvons inversement qu'une $f \in F \cap W$ est dans $TF$. Écrivons $f = pi$, avec $p \in TF$, $i \in TC$.
\[
\begin{tikzcd}[column sep=small, row sep=small]
X \arrow[rr, "f\in F\cap W"] \arrow[dr, "i\in TC"'] & & Y \\
& X' \arrow[ur, "p\in TF"'] &
\end{tikzcd}
\]
Comme $f \in F$, $i \in TC$, et $TC \longleftrightarrow F$ par définition de $F$, $\exists\, r : X' \longrightarrow X$ tel que $ri = \mathrm{id}_X$, $fr = p$. Donc $f$ est facteur direct de $p \in TF$, et comme $TF = C_{*}$ est stable par facteurs directs, on conclut que $f \in TF$.

\textbf{Corollaire 5.3.} \textit{Soit $(\mathcal{M}, W, C)$ une catégorie à cofibrations spéciale. Pour que $(W, C, F)$, $(F \defeq (TC)_{*})$ soit un triple de Quillen dans $\mathcal{M}$, il faut et il suffit que le couple $(TC, F)$ satisfasse à l'axiome de factorisation (i.e. que ce soit un précouple de Quillen). Pour que cette structure de Quillen soit clos, il faut et il suffit que $C$ soit stable par facteurs directs.} % [sic] soit clos -> soit close, agreement with "cette structure"

\msp{page 88}

\textsc{Démonstration}. La nécessité de l'axiome de factorisation pour $(TC, F)$, pour avoir un triple de Quillen, est tautologique. Pour la suffisance, ayant déjà les deux précouples de Quillen $(C, TF)$, $(TC, F)$ avec
% foot: Version 9 février 2019

\end{document}
